% Tapis de Quillen, with Quillen's offprint and preprint (cote 162-5, typed part): archive cote 162-5, site page 5 (PDF page 6).
% typescript. Draft compared: none.
% AI-transcribed by Claude Opus 5.5 (claude-opus-5-5, adaptive thinking, effort high), Message Batches API, batch t162_full (msgbatch_01UAh3zyWPd454AeDkhnmQ33), 2026-11-13; not yet checked by eye.
\scanpage{162-5/5}
\mspage{\hand{4}}

\xout{pleinementxfidèlex(respx} une équivalence). Même énoncé si on met dans le
coup un \xout{Anneauxlocx} système local d'anneaux sur X'. Enfin, f induit \xout{unx}
\xout{fonctepxpleinxfidxx(respx} une équivalence).
\xout{\unsure{isomorphismx}} entre la catégorie des coefficients locaux sur C et celle
\inserted[typed]{\xout{cettexpropriétéxdoit}}
des coéff. locaux sur C' \xout{(??xpeutxêtrexxlxxdoitxilxêtrexmisxdansxlesx}
\xout{hypothèsesxpréliminairesxsurxfx?)x} Principe de démonstration: on commence
par prouver ce dernier résultat, en notant que \xout{dansxlexcax} si un topos est
localement hom. trivial, il est loc\unsure{a} connexe et loc. simplement connexe,
d'où une bonne théorie du $\pi_0$ et du $\pi_1$ (qui sont ici discrets), et on est
ramené \xout{pour} à un cas particulier du critère d'homotopisme de Artin et
Mazur, savoir un critère \inserted[typed]{cohomologique} \unsure{(}pour qu'un homomorhisme de groupes $G\to H$ (ici de
groupes $\pi_1$ de X,X') soit un isomorphisme: il doit induire des isomorphis-%
mes sur les $H^{o}$ et $H^{1}$ (y inclus dans le cas non commutatif...). On prouve
la pleine fidélité en se ramenant par la suite spectrale aux cas des $\mathrm{Ext}^{\unsure{i}}$
de coéfficients locaux. Par la suite spectrale encore, cela résultera du
fait suivant: si X\deleted{'} est \unsure{localement} hom. trivial, alors la catégorie des
faisceaux abéliens loc. constants  est stable par $\underline{\mathrm{Ext}}^{i}$ , et le foncteur
$M\mapsto M_{X}$ \xout{desxfaisxxxxxxxconstants} (Ab) dans $X_{ab}$ commute aux dits $\mathrm{Ext}^{i}$ .
En fait, X et X' étant loc.hom\unsure{p}\hand{|}triviaux, les conditions suivantes sur f
seront équivalentes:

a) f est\xout{'} un homotopisme, i.e.  induit pour tout système local (pas
néc commutatif) sur X'un isomorphisme sur les $H^{i}$\xout{pourxixxm}.

b) \xout{Mêmexchosexave} f induit une équivalence $D^{b}_{lc}(X')\longrightarrow D^{b}_{lc}(X)$ .

a\hand{'}) f induit une équivalence entre la catégorie des systêmes locaux
abéliens sur X' et X, et des isomorphismes sur les $H^{i}$ correspondants (do
on ne prend ici que des coéfficients commutatifs).

\margin{?} J'ignore si on peut dans a) se borner aux systêmes locaux commutatif

\xout{ixexxdanx}
L'équivalence entre a) et b) fournit une première justification ou
motivation pour définir des types d'homotopie via la catégorie D\hand{$^{b}_{lc}$(X)},

\margin{5}
